\subsection{From Medical Records to Meaningful Insights}

To keep up with the evolution of digitalization, in the last 15-20 years more and more hospitals have opted for digital solutions when it comes to storing data and patient history. Thus, the classic paper records and files have been replaced with a new system (EHR - Electronic Health Records), which allows the digital storage of patient history. Because this data can be accessed much more easily than the classic ones, the analyses and reports made on patients or groups of patients have become easier to do. As shown in the review by Shen et al. \cite{shen2025}, the processing of this digital data can lead to valuable observations and analyses. In this way, common patterns can be discovered, treatments that have worked on other similar patients can be reapplied or the same recommendations can be applied to people in the same group. Therefore, the potential of this data is huge.

Electronic Health Records possess immense analytical value in the medical field, but the processing of this data is not always as simple as it seems at first glance. Patients may differ mathematically in small attributes, which can add noise to analyses and reports. Or simply medical data is private, under GDPR protection, and therefore difficult to use for academic research.

\subsubsection*{The First Barrier: Data Protection and Synthetic Data Accuracy}

The first major barrier encountered in the clinical analysis process is data protection, for data security reasons. Real EHR data contains sensitive information, such as patient names, dates of birth, or medical history. Due to GDPR restrictions, these are not accessible to the general public, requiring specific institutional approvals. Even if real data is obtained, sensitive fields must be anonymized due to the same GDPR rules.

The practical solution is to use synthetically generated data. This way, we can test the algorithms that will be used during the academic research without hitting privacy problems. We will use a dataset generated by Synthea. Although Synthea uses a probabilistic model to generate the data, its validity is supported by the principle of statistical fidelity. In their review \cite{walonoski2017}, Walonoski et al. demonstrate through prevalence graphs that as the simulated population grows, the results obtained become increasingly stable and closer to official statistics.

Mathematically, this process can be shown by the convergence of the synthetic distribution $P_{\text{synth}}$ to the real distribution $P_{\text{real}}$ as the measurement sample size $n$ grows:

\begin{equation}
    P_{\text{synth}} \xrightarrow{n \to \infty} P_{\text{real}}
\end{equation}

The ability to converge ensures that any data mining process conducted on this artificial database is theoretically generalizable to actual patients. Thus, it is safe to say that the application of synthetic EHRs does not undermine the validity of any data mining process.

\subsubsection*{The Second Barrier: Mathematical Complexity of Clinical Data}

The second major barrier occurs due to the variety of characteristics of the standalone data. The heterogeneity of the data makes it difficult to analyze the information using one approach. Specifically, clinical data of a certain patient consists of several kinds of events that are hard to accommodate by a single mathematical model. In particular, patient's clinical data may include:
\vspace{0.2em}
\begin{enumerate}
    \item \textbf{Data of different types:} a patient's record may belong to such types as integer (blood pressure, blood sugar, cholesterol), float (quality adjusted life year index, body mass index, hemoglobin) or string (patient's diagnoses, diagnoses codes, allergies).

    \item \textbf{Different data regarding variation:} some variables can differ in small intervals $[0.4, 4]$ - TSH or in larger intervals $[50, 500]$ - cholesterol. No single model will handle correctly these data, hence the necessity of their treatment before processing.

    \item \textbf{Interrupted or non-equally spaced timelines:} medical data may consist of interrupted records. A patient might have no information on his or her health condition for months, or even years (5 years without history - healthy patients), while other patients could have numerous interventions in a window of a few days.
\end{enumerate}
\vspace{1em}
Since there is no single approach to this problem, a solution should use a specific group  of algorithms and procedures to solve it optimally.

